When I was in high school, I could not wait for the physics I wanted to learn to be taught to me, so I taught myself. I worked through university-level theoretical mechanics, electromagnetism, relativity, quantum mechanics, and statistical mechanics on my own, and later advanced to the final training camp of the Physics Olympiad selection. The camp introduced me to far more advanced topics than I had met before, and I came home determined to become a theoretical physicist. Self-study also taught me something I use every day: how to find out what I do not know, and how to keep going when a text or a proof does not make sense.

That habit led me to teaching. In my first years at National Tsing Hua University I worked as a teaching assistant for Applied Mathematics, and since October 2022 I have been a tutor at the university's Counseling Center, where I support students with disabilities in their coursework. {\color{red} specific example (student) obstacle a student faced, what I tried, explained things.} This work taught me that explaining the same idea in one way does not work for everyone. I learned to listen first, to find out a student's obstacle, and to look for another route. It also affected how I think about my own field. If I can explain a concept to someone who approaches it very differently from me, then I understand it better myself. In September 2023 I received the university's Excellent Tutor Award for this work, and I count it among my most meaningful achievements, alongside my research.

Outside physics, I play the erhu, a traditional Chinese two-stringed instrument, in the National Tsing Hua University Chinese Orchestra. Playing in an ensemble taught me a kind of discipline that research does not. You practice alone, but the result only works when everyone listens to each other. In 2023 our group won first place in the national student chamber-music competition, and later that year we were invited to perform at the Hsinchu Chinese Music Festival. The experience as a member of an orchestra shapes how I collaborate with others, as I started to know .  {\color{red} continue}

%I hope to bring this tradition to UCSB, whether by performing or by sharing it with colleagues who have never heard it. {\color{red} what I hope?}

I grew up speaking Mandarin and learned to work in English through my coursework, seminars, and research. Because I have taught in one language and done research in another, I am used to moving between ways of explaining the same idea, which I expect will help me as a research or teaching assistant.

What I can offer the UCSB graduate community is a record of independent learning, a commitment to teaching, and experience supporting students who need a different approach. In graduate school I plan to keep teaching and mentoring, including by serving as a teaching assistant and by helping students who are new to the field. I also want to keep working on making physics learning accessible to students who learn differently. I hope to carry these habits, which I developed outside the lab as much as in it, into whatever research group I join.